\begin{lemma}[Finite energy and threshold counting]\label{thm:energy}
Let $P_0$ be a probability distribution on a finite set $Z$, and let $r_1,\ldots,r_K\in L^2(P_0)$. Set
\[
 \Gamma_{ij}=\E_{P_0}r_i r_j,\qquad \Lambda=\|\Gamma\|_\op.
\]
For every $a:Z\to\R$,
\begin{equation}\label{eq:energy}
 \sum_{j=1}^K\bigl(\E_{P_0}r_j a\bigr)^2
 \leq\Lambda\E_{P_0}a^2.
\end{equation}
In particular, if $|a|\leq1$ and $\tau>0$, then
\begin{equation}\label{eq:count}
 \bigl|\{j:|\E_{P_0}r_j a|>\tau\}\bigr|\tau^2\leq\Lambda.
\end{equation}
More generally, for any finite list of real numbers $c_i$, $\tau\geq0$, and $\sum_i c_i^2\leq\Lambda$, one has $|\{i:\tau<|c_i|\}|\tau^2\leq\Lambda$.
\end{lemma}
\begin{proof}
Let $T:\R^K\to L^2(P_0)$ send $b$ to $\sum_jb_jr_j$. Then $T^*T=\Gamma$. For every unit vector $b$,
\[
 |\langle b,T^*a\rangle|=|\langle Tb,a\rangle|
 \leq\|Tb\|_2\|a\|_2\leq\sqrt\Lambda\|a\|_2.
\]
Taking the supremum over $b$ and squaring proves \eqref{eq:energy}. Each selected summand in \eqref{eq:count} is greater than $\tau^2$; discarding nonnegative unselected summands proves the result. The same argument, with weak inequalities, proves the final statement even when $\tau=0$.
\end{proof}

\begin{theorem}[Joint-distribution SQ barrier]\label{thm:barrier}
Let $\mu$ be a full-support distribution on $X$ and let $P_0=\mu\otimes\mathrm{Unif}\{-1,+1\}$. For $j\in[K]$, let $P_j$ be the labeled distribution arising from some $D_j$ and $h_j\in\cH$. Put
\[
 r_j=\frac{P_j}{P_0}-1,\qquad
 \Gamma_{ij}=\E_{P_0}r_i r_j,\qquad
 \Lambda=\|\Gamma\|_\op,\qquad R=K/\Lambda.
\]
Any randomized adaptive $(m,\tau)$-SQ learner satisfying \eqref{eq:guarantee} obeys
\begin{equation}\label{eq:joint}
 1-2\epsilon\leq\sqrt{\frac\Lambda K}+
 \frac{m\Lambda}{K\tau^2},
 \quad\text{equivalently}\quad
 \frac m{\tau^2}\geq(1-2\epsilon)R-\sqrt R.
\end{equation}
Consequently, for $\epsilon<1/4$,
\begin{equation}\label{eq:Rbound}
 R\leq4+4m/\tau^2.
\end{equation}
\end{theorem}
\begin{proof}
First $\Lambda>0$. Indeed, with $a_j(x,y)=y h_j(x)$ we have $\|a_j\|_{L^2(P_0)}=1$ and
$\langle r_j,a_j\rangle=\E_{P_j}a_j-\E_{P_0}a_j=1$. Thus $\|r_j\|_2\geq1$ and $\Lambda\geq\Gamma_{jj}\geq1$.

For a fixed private seed, run the learner on the reference responses $\E_{P_0}q$. Write the resulting queries as $q_1,\ldots,q_m$ and the leaf output as $g_0$. This is a counterfactual tree path; the reference need not be realizable. For target $j$, choose the following memoryless oracle, which does not know the private seed:
\[
 O_j(q)=
 \begin{cases}
 \E_{P_0}q,&|\E_{P_j}q-\E_{P_0}q|\leq\tau,\\
 \E_{P_j}q,&\text{otherwise}.
 \end{cases}
\]
It is valid for $P_j$ on every query. If its run first departs from the reference transcript, that departure occurs at one of the reference queries with $|\langle r_j,q_t\rangle|>\tau$. Lemma~\ref{thm:energy} and a union bound give, for this fixed seed,
\[
 \frac1K\bigl|\{j:\text{the transcript departs from the reference}\}\bigr|
 \leq\frac{m\Lambda}{K\tau^2}.
\]
For the reference leaf, define $c_j=\E_{P_j}[y g_0(x)]$. Since $\E_{P_0}[y g_0(x)]=0$, Lemma~\ref{thm:energy} gives
\[
 \frac1K\sum_j|c_j|\leq
 \sqrt{\frac1K\sum_jc_j^2}\leq\sqrt{\Lambda/K}.
\]
Let $s_j$ be the correlation of the actual output under $O_j$. On a nondeparting transcript it equals $c_j$; on a departing transcript it is at most $1$. Therefore
$s_j\leq|c_j|+\1_{\text{departure for }j}$.
Average this inequality over $j$ and then over the private seed. The learner's guarantee gives an average correlation at least $1-2\epsilon$. This proves \eqref{eq:joint}. Notice that we did not choose a different oracle after seeing the seed.

Finally write $M=m/\tau^2$. From \eqref{eq:joint}, $R/2\leq\sqrt R+M$. The square $(\sqrt R-2)^2\geq0$ gives $\sqrt R\leq R/4+1$. Combining yields \eqref{eq:Rbound}.
\end{proof}

\paragraph{The requested fixed-marginal barrier.}
Set $D_j=\mu$ for every $j$. Then $r_j(x,y)=y h_j(x)$ and
\[
 \Gamma_{ij}=G_{ij}=\E_\mu[h_i h_j].
\]
For $R=K/\|G\|_\op\geq16$, \eqref{eq:joint} implies
\begin{equation}\label{eq:fixed}
 m\geq\tfrac14 R\tau^2.
\end{equation}
Indeed, $1-2\epsilon>1/2$ and $\sqrt R\leq R/4$. The cutoff matters: a uniformly positive lower bound proportional to $R\tau^2$ cannot hold for all tiny families, since a singleton needs no queries. The next theorem supplies the accompanying dimension lower bound.

\subsection{The retained sign-rank lower bound}
We use the spectral method associated with Forster \citep{Forster02}, but include its geometric step rather than treating the published theorem as a proof.

\begin{theorem}[Spectral sign-rank bounds]\label{thm:forster}
If $A\in\R^{N\times K}$ satisfies $|A_{xj}|\geq1$ for all entries, then
\begin{equation}\label{eq:forster}
 \sr(\sign A)\geq\frac{\sqrt{NK}}{\|A\|_\op}.
\end{equation}
If $M\in\{-1,+1\}^{N\times K}$ and $G=M^\top\diag(\mu)M$ for any distribution $\mu$ on its rows, then
\begin{equation}\label{eq:weighted}
 \sr(M)\geq\sqrt{K/\|G\|_\op}.
\end{equation}
\end{theorem}
\begin{proof}
Let $d=\sr(M)$, with $M=\sign A$ in the first assertion. Choose a strict representation
$M_{xj}=\sign\langle z_x,v_j\rangle$ in $\R^d$. If $d\geq K$, both conclusions are immediate: $\|A\|_\op\geq\sqrt N$ from any column, while $G_{jj}=1$ implies $\|G\|_\op\geq1$. We henceforth have $1\leq d<K$.

Perturb the finitely many $v_j$ within sufficiently small open neighborhoods so that signs are unchanged and every $d$ of them are independent. Such a perturbation exists: for each $d$-subset the determinant is a nonzero polynomial in the vector coordinates, and a finite union of its zero sets has empty interior. To justify this elementary fact, a nonzero univariate polynomial has finitely many roots; induction on the number of variables, by viewing the polynomial as a polynomial in the last variable, shows that a multivariate polynomial cannot vanish on a nonempty open box. Apply this to the product of the finitely many determinant polynomials.

We prove that there is an invertible $T$ such that the normalized vectors
$u_j=Tv_j/\|Tv_j\|$ satisfy
\begin{equation}\label{eq:isotropy}
 \sum_{j=1}^K u_ju_j^\top=(K/d)I_d.
\end{equation}
For every $1\leq r\leq d$, there is a common $\delta>0$ such that at most $d-r$ of the $v_j$ have projection norm less than $\delta$ onto any $r$-dimensional subspace. In fact, any $d-r+1$ of the vectors are independent, so cannot all belong to the orthogonal complement of an $r$-dimensional subspace. Their maximum projection norm is therefore positive for every such subspace. Its minimum is positive by compactness of the space of orthogonal projection matrices of rank $r$; this space is closed and bounded. Take the minimum over finitely many subsets and $r$.

For positive definite $Q$ with $\det Q=1$, minimize
\[
 F(Q)=\sum_{j=1}^K\log(v_j^\top Qv_j).
\]
Write the decreasing eigenvalues of $Q$ as $e^{s_1}\geq\cdots\geq e^{s_d}$, where $\sum_i s_i=0$. For each $r$, at least $K-d+r$ of the quadratic forms are at least $\delta^2e^{s_r}$. Order the $K$ log quadratic forms increasingly. For the first $d$ entries use respectively the bounds
$\log\delta^2+s_d,\ldots,\log\delta^2+s_1$; each of the remaining $K-d$ entries is at least $\log\delta^2+s_1$. Thus
\[
 F(Q)\geq K\log\delta^2+(K-d)s_1.
\]
Since $K>d$, every sublevel set has bounded largest eigenvalue. The determinant condition then bounds the smallest eigenvalue away from zero. Sublevel sets are compact, so a minimizer $Q$ exists.
For a symmetric traceless matrix $B$, the path
$Q(t)=Q^{1/2}\exp(tB)Q^{1/2}$ retains determinant one. Its derivative at zero gives
\[
 0=\sum_j\frac{v_j^\top Q^{1/2}BQ^{1/2}v_j}{v_j^\top Qv_j}
   =\tr\left(B\sum_j u_ju_j^\top\right),
 \qquad u_j=\frac{Q^{1/2}v_j}{\|Q^{1/2}v_j\|}.
\]
A symmetric matrix orthogonal to every traceless symmetric matrix is scalar: subtract its scalar trace part and take that difference as $B$. Taking traces gives \eqref{eq:isotropy}, with $T=Q^{1/2}$.

Replace each row vector by the normalization of $T^{-\top}z_x$, denoted $e_x$, so $\|e_x\|=1$ and all signs are preserved. Let $U$ be the matrix with columns $u_j$. For every row of $A$,
\[
 \langle e_x,UA_{x,\cdot}^{\top}\rangle
 =\sum_j |A_{xj}|\,|\langle e_x,u_j\rangle|
 \geq\sum_j\langle e_x,u_j\rangle^2=K/d.
\]
Here $|\langle e_x,u_j\rangle|\leq1$ was used. Consequently,
\[
 NK^2/d^2\leq\|UA^\top\|_F^2
 \leq\|A\|_\op^2\|U\|_F^2=K\|A\|_\op^2.
\]
This proves \eqref{eq:forster}. For $M$, use the same row inequality, square, and average under $\mu$:
\[
 K^2/d^2\leq\sum_x\mu(x)\|UM_{x,\cdot}^{\top}\|^2
 =\tr(UGU^\top)\leq\|G\|_\op\tr(UU^\top)=K\|G\|_\op.
\]
The last inequality follows by diagonalizing the positive semidefinite $G$, whose eigenvalues lie between zero and its operator norm. This proves \eqref{eq:weighted}.
\end{proof}

Theorems~\ref{thm:barrier} and \ref{thm:forster} prove the valid conditional assertion in the proposed spectral barrier: a certificate with ratio $R$ yields dimension lower bound $\sqrt R$ and query lower bound $\Omega(R\tau^2)$ when $R$ is large. They do not prove that all explicit sign-rank lower bounds admit such an ordinary label-Gram certificate. The following family disproves that extra assertion in precisely this form.

\subsection{Projective planes}
\begin{theorem}[An escaping plain-Gram family, but not an SQ separation]\label{thm:plane}
For each prime $q\geq2$, let $X_q$ be the points of the projective plane over $\mathbb F_q$, and let $\cH_q$ have one hypothesis per line, positive on that line and negative elsewhere. Set $N=q^2+q+1$ and $k=q+1$. Then:
\begin{align}
 \frac{q+1}{\sqrt q}&\leq\dc(\cH_q)\leq2q+3,\label{eq:planedc}\\
 \frac{|J|}{\|[\E_\mu h_i h_j]_{i,j\in J}\|_\op}&\leq16
 \quad\text{for every distribution $\mu$ and nonempty $J\subseteq\cH_q$}.
 \label{eq:planeplain}
\end{align}
There is a distribution-independent $N$-query tolerance-$\tau$ learner of loss at most $\tau$. For each \emph{fixed} $D$ and $0<\epsilon<1/4$, there is a learner allowed to depend on $D$, with tolerance $\epsilon/2$ and
\[
 m_D\leq1+s+\binom{s}{2},\qquad
 s=\left\lceil\frac{4\log N+2}{\epsilon}\right\rceil,
\]
that has loss at most $\epsilon$ for every target. Here $\log$ is natural.
Nevertheless every \emph{single} learner satisfying \eqref{eq:guarantee}, with $m\geq1$, $0<\tau\leq1$, and $\epsilon<1/4$, satisfies
\begin{equation}\label{eq:plane19}
 \dc(\cH_q)\leq19m/\tau^2.
\end{equation}
\end{theorem}
\begin{proof}
Points are one-dimensional subspaces of $\mathbb F_q^3$; lines are kernels of nonzero linear functionals, modulo nonzero scalar multiples. There are $(q^3-1)/(q-1)=N$ of each. A line contains $(q^2-1)/(q-1)=k$ points. Two distinct points span a unique two-dimensional subspace, hence lie on one common line. Two distinct kernels intersect in one one-dimensional subspace. These facts also show each point lies on $k$ lines.

Let $P$ be the $N\times N$ incidence matrix. Then
\[
 P\1=k\1,\qquad P^\top P=qI+\1\1^\top.
\]
Define $A=(N/k)P-\1\1^\top$. Its positive entries equal $N/k-1=q^2/(q+1)\geq1$, and its negative entries equal $-1$. Direct multiplication gives
\[
 A\1=0,\qquad
 A^\top A=(N/k)^2q\left(I-\frac{\1\1^\top}{N}\right).
\]
Indeed, the coefficient of $\1\1^\top$ before simplification is $N^2/k^2-N=-Nq/k^2$, since $k^2=N+q$. Thus $\|A\|_\op=(N/k)\sqrt q$, and Theorem~\ref{thm:forster} gives the lower bound in \eqref{eq:planedc}.

For the upper bound, enumerate the points by distinct integers $1,\ldots,N$ and use the common embedding
$\phi(t)=(1,t,\ldots,t^{2k})$. For a positive set $S$ of size $k$, the polynomial
\[
 p_S(t)=\tfrac12-\prod_{a\in S}(t-a)^2
\]
is positive on $S$ and negative on all other domain points: away from $S$ the product is a positive integer. Its coefficient vector is the required weight vector. Hence $\dc\leq2k+1$.

Now prove \eqref{eq:planeplain}. Write $K=|J|$, and restrict $P$ to the selected columns. If $K<16$, the Gram matrix has diagonal entries one, so its norm is at least one and the result follows. Suppose $K\geq16$ and let $r_x$ be the number of selected lines containing point $x$.
If $r_x<K/4$ for every $x$, take the unit vector $a=\1/\sqrt K$. At each point,
\[
 |\textstyle\sum_{j\in J}a_jh_j(x)|=|2r_x-K|/\sqrt K>\sqrt K/2.
\]
Its Gram Rayleigh quotient is at least $K/4$. Otherwise choose $x_0$ incident with $r\geq K/4$ selected lines and let $S$ be those $r$ lines. At $x_0$ all these labels are positive; at any other point at most one is positive, since two points have a unique common line. For the unit vector $a=\1_S/\sqrt r$,
\[
 |\textstyle\sum_j a_jh_j(x)|\geq(r-2)/\sqrt r\geq\sqrt r/2
 \quad\text{for all }x,
\]
since $r\geq4$. Its Rayleigh quotient is at least $r/4\geq K/16$, for every $\mu$. This proves \eqref{eq:planeplain}, including marginals concentrated on a single point.

The claimed learner queries $q_j(x,y)=y h_j(x)$ for every line and outputs a hypothesis with maximum answer. The true target has correlation one, hence receives answer at least $1-\tau$. A maximizing hypothesis has true correlation at least $1-2\tau$, and loss at most $\tau$, regardless of the marginal and the oracle. Thus it is one distribution-independent learner.

We prove the fixed-$D$ claim before excluding a common efficient learner. Let $\delta=\epsilon/2$ and sample a length-$s$ sequence from the known $D$. For any pair of hypotheses at $D$-distance greater than $\delta$, the probability of identical labels on the sequence is at most $(1-\delta)^s\leq e^{-\delta s}$. A union bound over fewer than $N^2$ pairs is at most $N^2e^{-\delta s}\leq e^{-1}<1$. Consequently some sequence has the property that hypotheses with identical traces are within $D$-distance $\delta$. Hardwire one such sequence and one representative hypothesis per trace into the learner. On at most $s$ distinct sampled points there are at most $1+s+\binom{s}{2}$ traces: the empty trace, at most $s$ singleton traces, and at most one line through each pair of distinct sampled points. Every trace of size at least two is determined by one such pair. The representatives therefore form a $D$-cover of the claimed size. Query their label correlations and select the largest answer. A representative within $\delta$ of the true target has correlation at least $1-2\delta$. With tolerance $\tau=\epsilon/2$, the selected output has correlation at least $1-2\delta-2\tau$ and error at most $\delta+\tau=\epsilon$. The hardwired cover depends on $D$; this is not one learner serving all marginals.

To rule out a much cheaper common learner, choose for target line $j$ the distribution
\[
 D_j(x)=\begin{cases}(2k)^{-1},&P_{xj}=1,\\[1mm]
 (2(N-k))^{-1},&P_{xj}=0.
 \end{cases}
\]
Take the reference $P_0$ to be uniform on $X_q\times\{-1,+1\}$. With $c=k/(N-k)$, its residual likelihood ratios satisfy
\[
 r_j(x,+1)=A_{xj},\qquad r_j(x,-1)=-cA_{xj}.
\]
For example the first formula is $2N D_j(x)-1=N/k-1$ on the line and $-1$ elsewhere; the second is $-1$ on the line and $N/(N-k)-1=c$ elsewhere. Consequently
\begin{align*}
 \Gamma&=\frac{1+c^2}{2N}A^\top A,\\
 \Lambda&=\frac{qN}{2}\left(\frac1{(q+1)^2}+\frac1{q^4}\right),\\
 R&=\frac{2}{q((q+1)^{-2}+q^{-4})}\geq\frac{(q+1)^2}{q}\geq q.
\end{align*}
The penultimate inequality uses $q^4\geq(q+1)^2$ for $q\geq2$, which follows from $q^2\geq q+1$. By Theorem~\ref{thm:barrier}, $R\leq4+4M$ for $M=m/\tau^2$. Therefore
\[
 \dc(\cH_q)\leq2q+3\leq2R+3\leq11+8M\leq19M,
\]
where the final inequality uses $M\geq1$.
\end{proof}

\paragraph{Why projective planes do not separate.}
The family has unbounded dimension while all the proposed ordinary label-Gram ratios remain bounded. This is a genuine failure of the stated universality premise, not a counterexample to OQ2. The efficient learner needed for such a separation is impossible here: the target-dependent balanced marginals produce a joint-distribution ratio of order $q$. The lower-bound witness changes entry magnitudes, which the ordinary label Gram matrix cannot capture. This also verifies an explicit separation between fixed-distribution and distribution-independent SQ learning: $m_D=O(n^2/\epsilon^2)$ at tolerance $\epsilon/2$, whereas for $q\geq16$ any common learner needs $m/\tau^2\geq q/4$. Since $N\leq2q^2$, the latter is at least $2^{(n-1)/2}/4$. The dimension lower bound is at least $2^{(n-1)/4}$. None of these fixed-$D$ learners can be substituted for the single algorithm in OQ2.

\subsection{Transcript compression}
\begin{theorem}[Marginal-branching and approximate probabilistic bounds]\label{thm:tree}
Let $B=\lceil2/\tau\rceil+1$, with $0<\tau\leq1$. For a deterministic learner satisfying \eqref{eq:guarantee}, decompose each query as
\[
 q(x,y)=a(x)+yb(x),\quad
 a(x)=\tfrac12(q(x,+1)+q(x,-1)),\quad
 b(x)=\tfrac12(q(x,+1)-q(x,-1)).
\]
Construct its marginal-response tree as follows. If $a$ is constant, retain only the answer equal to that constant. Otherwise retain the $B$ equally spaced answers in $[-1,1]$. Let $u$ be the maximum number of nonconstant-$a$ nodes on a retained root-to-leaf path. Then
\begin{equation}\label{eq:tree}
 \dc(\cH)\leq(m+1)B^u.
\end{equation}
In particular, a deterministic correlational-only learner has $\dc(\cH)\leq m+1$.
For an arbitrary randomized learner satisfying \eqref{eq:guarantee},
\begin{equation}\label{eq:approx}
 \dc^{\epsilon}(\cH)\leq B^m.
\end{equation}
\end{theorem}
\begin{proof}
The grid has spacing $2/(B-1)\leq\tau$, so nearest-grid rounding errs by at most $\tau/2$. The marginal tree has at most $B^u$ leaves: suppress unary nodes and use induction on the number of remaining branching nodes per path. It has at most $(m+1)B^u$ nodes, because every node has a descendant leaf and a depth between zero and $m$, and assigning each node one such pair is injective when the depth is fixed along a leaf's path.
Collect the function $b$ at every internal node and the Boolean output at every leaf. These are the common features $f_1,\ldots,f_s$, where $s\leq(m+1)B^u$ and each feature has absolute value at most one.

Fix $D,h$. Follow the tree by giving the exact constant response or rounding $\E_Da$ to the grid. If at some node $|\E_D hb|>\tau/2$, one of the signed features $b,-b$ has correlation greater than $\tau/2$. Otherwise every answer is a valid SQ response: the difference from the true expectation $\E_Da+\E_D hb$ is at most $\tau/2+\tau/2$. A valid policy realizing this finite transcript extends to all other histories by exact answers. The deterministic guarantee then gives leaf correlation at least $1-2\epsilon>1/2\geq\tau/2$. We have proved
\begin{equation}\label{eq:minimaxpremise}
 \forall D,\qquad \max_{i,\ s_i\in\{-1,+1\}}\E_D[h(x)s_if_i(x)]\geq\tau/2.
\end{equation}
We supply the finite separation step. Let $C$ be the convex hull of the vectors $(h(x)s_if_i(x))_{x\in X}$ and let $Q=(\tau/2)\1+\R_{\geq0}^X$. If $C\cap Q$ were empty, the compact set $C$ and closed set $Q$ would have a positive attained minimum distance. To see attainment, take a minimizing sequence; the $C$ components have a convergent subsequence, and bounded distances make the $Q$ components bounded too. A zero minimum would put the limit in both sets. Let $(c,z)$ minimize the positive distance and put $v=z-c$. By differentiating squared distance along line segments, for every $c'\in C,z'\in Q$ we obtain
\[
 \langle v,c'-c\rangle\leq0,\qquad
 \langle v,z'-z\rangle\geq0.
\]
Taking $z'=z+te_x$ gives $v_x\geq0$. Furthermore $v\ne0$, and
\[
 \max_{c'\in C}\langle v,c'\rangle
 \leq\langle v,c\rangle<\langle v,z\rangle
 =\min_{z'\in Q}\langle v,z'\rangle=(\tau/2)\sum_xv_x.
\]
Normalizing $v$ to a distribution contradicts \eqref{eq:minimaxpremise}. Hence a convex combination in $C\cap Q$ exists. Its coefficients define a linear combination of the common features whose product with $h(x)$ is at least $\tau/2$ at every point. This is a strict representation and proves \eqref{eq:tree}. For correlational-only queries $a=0$, hence $u=0$.

For a randomized learner, fix a seed and instead retain all grid answers to each full query. There are at most $B^m$ leaves. Use all leaf outputs as coordinates of an embedding, padding when needed. This distribution over embeddings depends only on the learner's seed, not on $D,h$. For fixed $D,h$, the policy rounding each true expectation to the grid is valid. Its output is one of the leaf features. Selecting that coordinate as $w$ gives
\[
 \inf_w L_{D,h}(\sign\langle w,\phi_\omega\rangle)
 \leq L_{D,h}(A_\omega^{O_{D,h}}).
\]
Averaging over the seed proves \eqref{eq:approx}. No seed is selected after observing $D$ or $h$.
\end{proof}

The deterministic argument does not derandomize the learner: fixing its seed need not preserve its expected guarantee. The probabilistic bound is exponential, not polynomial, and gives no exact $\dc_{1/2}$ bound.

